\scanpage{24}
a) Le foncteur $\unsure{w|\omega}_N$ (de la catégorie des $k$-modules vers celle des $S$-modules gradués) est pleinement fidèle.

b) \inserted{$\forall M$,} l'application
\[
\deleted{\unsure{x}}M\xrightarrow{\varphi}\prod_{i\geq0}w_N(M)_i=\prod_{i\geq N}M.T^{(i)}
\]
donnée par
\[
\varphi(x)=(x.T^{(i)})_{i\geq N}
\]
induit une bijection de $M$ avec le sous-ens. de $\prod_{i\geq N}\deleted{\unsure{M}}w_N(M)_i$ formé des systèmes $(x_i)_{i\geq N}$ satisfaisant la condition
\[
T^{(j)}x_i=c_{ij}x_{i+j}\qquad\forall i\geq N,\ j\geq1\tag*{$(*)$}
\]

\uline{Dém}

\deleted{a)}b) $x_i=xT^{(i)}$, on a $T^{(j)}x_i=xT^{(i)}T^{(j)}=c_{ij}xT^{(i+j)}=c_{ij}x_{i+j}$, donc $\varphi(x)$ se trouve dans le sous-ens indiqué. L'injectivité de $\varphi$ est claire. Bijectivité ? Posons $x_i=\xi_i T^{(i)}$, $\xi_i\in M$. La condition $(*)$ s'écrit
\[
\begin{aligned}
\xi_i T^{(i)}T^{(j)}&=c_{ij}\xi_{i+j}T^{(i+j)}\\[-2pt]
\shortparallel&\\[-2pt]
c_{ij}\xi_iT^{(i+j)}&
\end{aligned}
\]
i.e.
\[
c_{ij}(\xi_{i+j}-\xi_i)=0\qquad\text{pour}\quad i\geq N,\ j\geq1,
\]
et il résulte du th. que les $\xi_i$ sont tous égaux à un $=$ $\xi\in M$, cqfd.

a) \unsure{Soient|Soit} $M,M'$ sur $k$. Alors $\operatorname{Hom}(w_N(M),w_N(M'))$ s'identifie à l'ens des systèmes de $k$-hom $(i\geq N)$
\diagram{\begin{tikzcd}[ampersand replacement=\&,column sep=large,row sep=small]
u_i\colon M \arrow[r,phantom,"\unsure{-|\to}"] \arrow[d,phantom,"\rotatebox{90}{$\simeq$}"] \& M' \arrow[d,phantom,"\rotatebox{90}{$\simeq$}"]\\
M.T^{(i)} \arrow[d,equal] \& M'.T^{(i)} \arrow[d,equal]\\
w_N(M)_i \& w_N(M')_i
\end{tikzcd}}
% archive sheet number: 23

